\documentclass{article}
\usepackage{geometry}
\usepackage{concmath}
\usepackage[OT1]{fontenc}
\usepackage{graphicx}
\usepackage{amsfonts}
\usepackage{amsmath,mathtools}
\usepackage{calrsfs}
\usepackage{color}
\usepackage[dvipsnames]{xcolor}
\usepackage[many]{tcolorbox}
\usepackage{setspace}
\usepackage{tkz-tab}    
\usepackage{amssymb}          
\usepackage{multicol}  
\usepackage{esint}
\geometry{margin=0.5in}
\definecolor{green}{RGB}{52, 235, 168}
\definecolor{violet}{RGB}{165, 83, 237}
\definecolor{lemon}{RGB}{170, 237, 83}
\definecolor{red}{RGB}{240, 125, 84}
\definecolor{blue}{RGB}{105, 255, 242}

\begin{document}
  
\noindent \colorbox{green}{541.1}\\
\begin{equation*}
    I = \int_{0}^{+\infty}\frac{t^{3}}{e^{t}-1}dt 
\end{equation*}
* \(t \mapsto \frac{t^{3}}{e^{t}-1}\) continue sur \(]0, +\infty[\)\\
* Au voisinnage de \(t \rightarrow +\infty\): 
\begin{align*}
  \frac{t^{3}}{e^{t}-1} \sim t^{3}e^{-t} &= o(e^{-t})\\
                                         &\hookleftarrow \text{ intégrable en } +\infty
\end{align*}
* Au voisinnage de \(t \rightarrow 0\): 
\begin{align*}
  \frac{t^{3}}{e^{t}-1} &= \frac{t^{3}}{1 + t - 1 + o(t)} \\
    &\sim \frac{t^{3}}{t} \\
    &= t^{2}
\end{align*}
\begin{align*}
  I &= \int_{0}^{+\infty}t^{3}e^{-t}\times \frac{1}{1-e^{-t}}dt \\
    &= \int_{0}^{+\infty}t^{3}e^{-t}\sum_{k=0}^{+\infty}(e^{-t})^{k}dt\\
    &= \int_{0}^{+\infty}\sum_{k=1}^{+\infty}t^{3}(e^{-t})^{k}dt
\end{align*}
\(\forall k, t, t^{3}(e^{-t})^{k} \geq 0\)\\
On calcule dans \([0, +\infty[\cup\{+\infty\}\)

\noindent \colorbox{green}{541.2}\\
\begin{align*}
  \forall t > 0. \frac{t^{3}}{e^{t}-1}&= t^{3}\frac{e^{-t}}{1-e^{-t}}\\ 
                                      &= t^{3}\sum_{n=1}^{+\infty}(e^{-t})^{n}\\
                                      &= \sum_{n=1}^{+\infty}\underbrace{t^{3}e^{-nt}}_{f_{n}(t)}\\ 
  S_{N}(t) &= \sum_{n=1}^{N}t^{3}e^{-nt}\\ 
  \int_{0}^{+\infty}S_{N}(t) &= \int_{0}^{+\infty}\sum_{n=1}^{N}f_{n}(t)dt \\ 
                             &= \sum_{n=1}^{N}\int_{0}^{+\infty}f_{n}(t)dt \text{ par linéarité}
\end{align*}
avec \(\int_{0}^{+\infty}t^{3}e^{-nt}\)\\
Par partie sous réserve de limite finie du crochet et convergence 
d'une intégrale 
\begin{align*}
  \int_{0}^{+\infty} t^{3}e^{-nt} &= \frac{6}{n^{4}}\\
  \int_{0}^{+\infty}S_{N}(t)  &= \sum_{n=1}^{+\infty}\frac{6}{n^{4}}
\end{align*}
Convergence simple : 
\begin{align*}
  S_{N}(t)\xrightarrow[N \rightarrow +\infty]{} S(t)
\end{align*}
\underline{Dominons}: 
\begin{align*}
  |S_{N}(t)| &= |\sum_{n=1}^{N}t^{3}e^{-nt}|\\
             &\leq \sum_{n=1}^{+\infty}t^{3}e^{-nt}\\ 
             &= \underbrace{\frac{t^{3}}{e^{t}-1}}_{\varphi(t)} \text{ indépendant de N}
\end{align*}
* \(\varphi(t)\) continue sur \(]0, +\infty[\)\\
* \underline{Au voisinnage de \(t \to 0\)}\\
\begin{align*}
  \varphi(t) &= \frac{t^{3}}{1+t+o(t)}\\
             &\backsim t^{2}
\end{align*}
* \underline{Au voisinnage de \(t \rightarrow +\infty\)}\\
\begin{align*}
  \varphi(t) &= \frac{t^{3}}{e^{t}-1}\\ 
             &= \frac{o(e^{\frac{t}{2}})}{e^{t}}\\
             &= o(e^{-\frac{t}{2}})
\end{align*}

\noindent \colorbox{green}{541.14}\\
\begin{align*}
  \int_{0}^{1}\frac{1}{x^{x}}dx &= \int_{0}^{1}e^{-x\ln x}dx \\
                                &= \int_{0}^{1}\sum_{n=0}^{+\infty} \underbrace{\frac{(-x\ln x)^{n}}{n!}}_{f_{n}(x)}dx
\end{align*}
* \(\sum_{}^{}f_{n}\) \\
Donc \(\int_{0}^{1}\sum_{n=0}^{+\infty}f_{n}(t) dt = \sum_{n=0}^{+\infty} f_{n}(t)dt\)
dans  \([0, +\infty[ \cup \{+\infty\}\)\\
\begin{align*}
  \int_{0}^{1}(-x\ln(x))^{n}dx &= \int_{0}^{1}x^{n} (-\ln(x))^{n}dx \\
                               &= [\frac{x^{n+1}}{n+1}- (-\ln(x))^{n}]^{1}_{0} - \int_{0}^{1}\frac{x^{n+1}}{n+1}.(-)\frac{n}{x}(-\ln(x))^{n-1}dx\\ 
                               &= 0 - 0 + \frac{n}{n+1}\int_{0}^{1}x^{n}(-\ln(x))^{n-1}dx \\
                               &= \frac{n}{n+1}\frac{n-1}{n}\int_{0}^{1}x^{n}(-\ln(x))^{n-2}dx \\
                               &= \frac{n.(n-1)\dots}{(n+1)(n+1)\dots}\frac{1}{n+1}\int_{0}^{1}x^{n}(-\ln(x))^{0}dx\\
                               &= \frac{n!}{(n+1)^{n+1}}
\end{align*}
\begin{align*}
  \int_{0}^{1}\frac{1}{x^{x}}dx &= \sum_{n=0}^{+\infty}\frac{1}{(n+1)^{n+1}}\\
                                &= \sum_{n=1}^{+\infty}\frac{1}{n^{n}} < +\infty
\end{align*}
Donc les intégrales et séries sont convergentes \\

\noindent \colorbox{green}{541.12}\\
\begin{equation*}
    f_{n}(x) = x^{n}(1-\sqrt{x})
\end{equation*}
a/ Soit \(x \in \mathbb{R}\) fixé\\ 
\begin{equation*}
    \sum_{}^{}x^{n}(1-\sqrt{x})
\end{equation*}
* \underline{Si \(x < 1\):}\\
\begin{align*}
  x^{n}(1-\sqrt{x})=O(x^{n}) &\text{ tg série convergente}\\
  \sum_{n=0}^{+\infty}x^{n}(1-\sqrt{x})&= |1-\sqrt{x}|\frac{1}{1-x}\\
                                       &= \frac{1-\sqrt{x}}{1-x}
\end{align*}
* \underline{Si \(x > 1\)}\\
\(\sum_{n=0}^{+\infty}x^{n}(1-\sqrt{x})\) diverge grossièrement donc 
\(\sum_{}^{}f_{n}\) diverge\\

\noindent \(\sum_{}^{}f_{n}\) converge simplement sur \(]0, 1]\) 

\noindent b/ Pour \(x \in [0, 1]\)\\
\underline{Si x = 1}
\begin{align*}
    \sum_{n=0}^{+\infty}f_{n}(x) = 0
\end{align*}
\underline{Si \(x < 1\)}
\begin{align*}
  \sum_{n=1}^{+\infty}f_{n}(x) &= (1-\sqrt{x})(\frac{x}{1-x})\\
                               &= \frac{x}{1+\sqrt{x}}
\end{align*}
\noindent
c/ On pose \(f : x \mapsto 
\begin{cases}
  0 \text{ si x = 1}\\ 
  \frac{x}{1+\sqrt{x}} \text{ si } x < 1
\end{cases}
\)\\
Par non transfert de continuité, \(\sum_{}^{}f_{n}\) ne converge
pas uniformément vers \([0, 1]\)\\ 

\noindent
d/
\begin{align*}
  R_{n}&= \sum_{k = n+1}^{+\infty}(1-\sqrt{x})x^{k}\\
       &= (1-\sqrt{x})x^{n+1}\frac{1}{1-x}\\
       &= \frac{x^{n+1}}{1+\sqrt{x}}\\
  \int_{0}^{1}\frac{x^{n+1}}{1+\sqrt{x}}dx &\text{ Or \(\frac{x^{n+1}}{1+\sqrt{x}}\xrightarrow[n\to+\infty]{} 0\) sur ]0, 1[}\\ 
  R_{n}&\xrightarrow[n \to +\infty]{cs} x\mapsto 0
\end{align*}
\begin{align*}
  R_{n}(x) &= \frac{x^{n+1}}{1+\sqrt{x}}\\
           &\leq x^{n+1}\\
           &\leq 1 \text{ intégrable sur \(]0, 1[\)}
\end{align*}
Par convergence dominée, 
\begin{align*}
  \Aboxed{\int_{0}^{1}R_{n}(x)dx \xrightarrow[n \rightarrow +\infty]{} 0} 
\end{align*}\\

\noindent 
e/ 
\begin{align*}
  \sum_{n=1}^{+\infty} f_{n}(x) &= \sum_{n=1}^{N}f_{n}(x) + R_{N}(x)\\
  \int_{0}^{1}\sum_{n=1}^{+\infty}f_{n}(x)dx &= \int_{0}^{1}\sum_{n=1}^{N}f_{n}(x)dx + \int_{0}^{1}R_{N}(x)dx\\
                                             &= \sum_{n=1}^{N}\int_{0}^{1}f_{n}(x)dx + \underbrace{\int_{0}^{1} R_{N}(x)dx}_{\rightarrow 0}
\end{align*}
Donc \(\sum_{n=1}^{+\infty}\int_{0}^{1}f_{n}(x)dx\) converge et \\
\begin{align*}
    \Aboxed{\int_{0}^{1}\sum_{n=1}^{+\infty}f_{n}(x)dx = \sum_{n=1}^{+\infty}\int_{0}^{1}f_{n}(x)dx}
\end{align*}

\noindent 
f/ 
\begin{align*}
  \int_{0}^{1}f_{n}(x)dx &= \int_{0}^{1}x^{n}(1-\sqrt{x})dx\\ 
                         &= \int_{0}^{1}x^{n}dx - \int_{0}^{1}x^{n+\frac{1}{2}}dx\\
                         &= [\frac{x^{n+1}}{n+1}]^{1}_{0}- [\frac{x^{n+\frac{3}{2}}}{n+\frac{3}{2}}]_{0}^{1}\\ 
                         &= \frac{1}{n+1}-\frac{1}{n+\frac{3}{2}}\\
                         &= \frac{1}{n+1}-\frac{2}{2n+3}\\
                         &= \frac{1}{(n+1)(2n+3)}
\end{align*}
D'autre part, on pose \(x=t^{2}, dx = 2t dt\)
\begin{align*}
  \int_{0}^{1}\frac{x}{1+\sqrt{x}}dx &= 2 \int_{0}^{1}\frac{t^{3}}{1+t}dt \\
                                     &= 2 \int_{0}^{1} t^{2} - t + 1 \frac{1}{1+t}dt \\
                                     &= 2[\frac{1}{3}-\frac{1}{2}+1 - \ln 2]
\end{align*}














\end{document}

